\section{Proofs of the discrete limit distribution and the Skellam rule}\label{app:inversion}
We prove Theorems~\ref{thm:cross}, \ref{thm:uniform}, and \ref{thm:coverage}. Write the iid coding coefficients as $\varepsilon_i$. By \eqref{eq:raw-expand}, $\widetilde d-d=\frac2k\sum_{i<j}H_{ij}\varepsilon_i\varepsilon_j$ holds for every sample. The final subsection discusses an extension with multiple sign counts per cell.

\subsection{Exact moments for a specified duplicate edge pair}
Let $S_1=S_2$, with the remaining $d-2$ edges independent and uniform. If the two coding coefficients have opposite signs, the edges cancel completely, leaving the remaining $d-2$ elements. If they are equal, the edges combine into a single edge with coefficient $\pm2$. For $d\ge2$,
\begin{equation*}
\E[\widehat d\mid S_1=S_2]=d,\qquad
\Var(\widehat d\mid S_1=S_2)=\frac{2d(d-1)}{M-1}+4-\frac4{M-1}.
\end{equation*}
The component with opposite coefficients has variance $2(d-2)(d-3)/(M-1)$, and the component with equal coefficients has variance $2(d-2)(d+5)/(M-1)$. Averaging these variances and adding $4$ from the difference between the component means $d\mp2$ gives the result.

\subsection{From failure conditioning to the planted distribution}
Lemma~\ref{lem:plant} gives $\|\mathcal L(\mathcal G\mid F)-Q\|_{TV}\to0$ for all $2\le d=o(M)$. Here $Q$ selects a pair of edge labels uniformly and then conditions on the two edges having the same mapping. For $2\le d\le C\sqrt M$, both $\lambda$ and $d^2/M^k$ are $O_k(M^{1-k})$. Since $b_k=\lceil3(k-2)/2\rceil$, we have $k-b_k-3=-\lceil k/2\rceil\le-2$ and hence $dM^{k-b_k-3}=O_k(d/M^2)$. The distance is therefore at most $O_k(d/M^2+M^{1-k})$. Adding independent signs and applying \eqref{eq:raw-expand} preserve this TV bound. The statistic is invariant under label permutation, so we may fix the duplicate edge pair as 1, 2. Their common edge and the remaining $d-2$ edges are independent and uniform. For $d=2$, the null distribution is exact under the planted distribution. We treat empty sums below as zero.

\subsection{Three negligible event classes for bounded \texorpdfstring{$c$}{c}}
By the union bound, the common edge of the duplicate edge pair intersects a remaining edge with probability $O_k(d/M)$. Any specified pair of remaining edges shares at least two cells with probability $O_k(M^{-2})$. Summing over all such pairs gives $O_k(d^2/M^2)$. In the intersection graph of the remaining edges, for any three specified edges, the probability that one specified edge intersects both others is
\begin{equation*}
\Bigl(1-\binom{M-k}{k}\Big/\binom Mk\Bigr)^2=O_k(M^{-2}),
\end{equation*}
This equality follows by conditioning on the edge shared by the two intersection tests. The other two edges are independent, and their intersection probability with the shared edge is independent of its position. Summing over ordered triples bounds the probability that the intersection graph is not a matching by $O_k(d^3/M^2)$. Thus, outside an event with probability bounded by \eqref{eq:tv-error}, the common edge of the duplicate edge pair is disjoint from the remaining edges, and all overlaps among the remaining edges are single-cell overlaps forming a matching.

\subsection{Poisson approximation for overlaps among the remaining edges}
For each pair of remaining edges, consider the indicator that they intersect. These indicators have a dependency graph defined by shared edge labels. Collections of indicators with disjoint label sets are independent. Each indicator has mean $p=1-\binom{M-k}k/\binom Mk=k^2/M+O_k(M^{-2})$ and $O(d)$ neighbors. By the calculation in the previous subsection, two indicators sharing an edge have joint probability exactly $p^2$. We apply the local dependence bound of Arratia, Goldstein, and Gordon \cite[Theorem~1]{arratia1989chen}. The neighborhood includes the indicator itself, and indicators outside it depend only on the remaining independent edges, so the third term is zero. The bound is a universal constant times the sum of products of means within neighborhoods plus the sum of joint probabilities over distinct neighbors. Both sums are $O_k((d^2+d^3)/M^2)$ here. Thus the number of intersecting edge pairs among the remaining edges has TV distance $O_k((d^2+d^3)/M^2)$ from $\Pois(\binom{d-2}2p)$. We replace the parameter by $\binom{d-2}2k^2/M=2\lambda_d$. The TV distance between two Poisson distributions is at most the difference between their parameters: we couple the larger-parameter variable to the smaller one by adding an independent Poisson increment, and the variables differ only when this increment is nonzero. The replacement adds an error of $O_k(d^2/M^2)$.

\subsection{Independence of the overlap sign marks and the Skellam limit}
On the matching event, each overlap uses a pair of iid signs shared with no other overlap. Their products are therefore mutually independent and symmetric, and independent of the duplicate edge pair's sign product $\varepsilon_1\varepsilon_2$. Given the mapping, $\widetilde d$ depends only on the number of overlaps among the remaining edges and these independent $\pm1$ marks. To formalize the coupling, we retain the true overlap count and generate that many new independent marks for each mapping. On the matching event, the old and new marked sums have the same conditional distribution. On all other mappings, we couple them arbitrarily, with a distributional error at most the exceptional event's probability. We then apply the same marking kernel to the Poisson approximation of the count, preserving the TV bound. By Poisson thinning \cite{kingman1993poisson}, $\Pois(2\lambda_d)$ independent $\pm1$ marks split by sign into two independent $\Pois(\lambda_d)$ counts. Their difference, $Z_d$, has a Skellam distribution \cite{skellam1946}. Adding the duplicate edge pair's contribution $2\varepsilon_1\varepsilon_2$ gives $J_{\lambda_d}$ in \eqref{eq:J}. Thus $\widetilde d$ is close to the null distribution $d+J_{\lambda_d}$ in \eqref{eq:aux}. We use the Poisson approximation for the unconditional count and the matching event to bound the coupling error.

The four steps give a total error of $O_k(d/M+d^2/M^2+d^3/M^2+d/M^2+M^{1-k})$. Since $d\le M$ and $k\ge3$, the terms $d/M$ and $d^3/M^2$ absorb all others, giving \eqref{eq:tv-error}. When $d/\sqrt M\to c$, we have $\lambda_d\to k^2c^2/4$. Continuity in the Poisson parameter gives convergence of $J_{\lambda_d}$ in total variation to $J_{k^2c^2/4}$, proving Theorem~\ref{thm:cross}(ii). The $\pm2$ and Skellam terms are independent, so $\Var J_\lambda=4+8\lambda/k^2$ and the limiting variance is $4+2c^2$. For (iii), as $\lambda\to0$, we have $\Pr(Z\ne0)\le2\lambda\to0$ and $J_\lambda\Rightarrow2\varepsilon_1\varepsilon_2$. As $\lambda\to\infty$, we have $Z/\sqrt{2\lambda}\Rightarrow\mathcal N(0,1)$ and $\sqrt{8\lambda}/k=\sqrt2\,c$. The $\pm2$ term divided by $\sqrt2\,c$ tends to zero, so $J_{k^2c^2/4}/(\sqrt2\,c)\Rightarrow\mathcal N(0,1)$. The estimator $\widehat d$ has the same weak limit: the AMS term in \eqref{eq:raw-decomp} has coefficient $k/M$, and $(\sum_i\varepsilon_i)^2=O_p(d)=O_p(\sqrt M)$. \qed

\subsection{Sparse limit for fixed small \texorpdfstring{$d$}{d}}
\begin{proposition}[Limit when $d=o(\sqrt M)$]\label{prop:smalld}
Fix $k\ge3$ and let $2\le d=o(\sqrt M)$. Conditional on failure, both $\widehat d-d$ and $\widetilde d-d$ converge weakly to the distribution assigning equal probability to $\pm2$. The Gaussian rule's conditional miss rate tends to zero for $d=2,3$, and to $1/2$ for fixed $d\ge4$ or $d\to\infty$.
\end{proposition}

\begin{proof}
Theorem~\ref{thm:cross} gives the equally weighted $\pm2$ limit for the fluctuation of $\widetilde d$ when $d=o(\sqrt M)$. By \eqref{eq:raw-decomp}, the correction for $\widehat d$ is $O_p(d^2/M)=o_p(1)$, so its fluctuation has the same weak limit. Consider the event that the remaining edges are pairwise disjoint and disjoint from the duplicate edge pair. On this event, $\widetilde d$ equals $d+2\varepsilon_1\varepsilon_2$ exactly, and the event has conditional probability tending to one. The Gaussian rule has $t(d)=d-O(s(d))$ with $s(d)\to0$. The negative component leads to rejection of the true candidate, and the positive component leads to acceptance. For $d=2,3$, the negative component gives an observation below $t(2)$, so the rule returns $+\infty$, which is not a miss. For $d\ge4$, the negative component still leads to acceptance of candidate $d-2$, causing a miss. For $d=2$, failure is equivalent to a duplicate edge pair, and $\Pr(\widehat d=0\mid F)=\Pr(\widehat d=4\mid F)=1/2$ for every $M>k$. Thus multiplying $\widehat d$ by any fixed constant to obtain an upper bound leaves a post-failure risk of at least $1/2$.
\end{proof}

\subsection{Proof of the uniform CDF approximation}
\begin{proof}
First consider $d/\sqrt M\to\infty$. Appendix~\ref{app:clt} gives a conditional Gaussian limit for $\widetilde d$ standardized by $s(d)$. The null distribution is the distribution of $d+J_{\lambda_d}$. As $\lambda_d\to\infty$, we can express $Z_d$ in the Skellam term as a sum of iid increments whose number grows with $\lambda_d$. The CLT gives $(2/k)Z_d/\sqrt{8\lambda_d/k^2}\Rightarrow\mathcal N(0,1)$. The ratio of $8\lambda_d/k^2=2(d-2)(d-3)/M$ to $s(d)^2$ tends to one, and the $\pm2$ term divided by $s(d)$ tends to zero. Hence $J_{\lambda_d}/s(d)\Rightarrow\mathcal N(0,1)$. Both standardized CDFs converge to the continuous CDF $\Phi$. Weak convergence to a continuous limit implies uniform CDF convergence (P\'olya's theorem), so their Kolmogorov distance tends to zero.

To prove uniformity, suppose that there exist $\delta>0$ and a sequence $(M,d)$ along which the distance is at least $\delta$. Take a subsequence with $d/\sqrt M\to c\in[0,\infty]$. If $c<\infty$, this subsequence eventually satisfies $d\le C\sqrt M$ for some constant, and Theorem~\ref{thm:cross} gives a TV distance tending to zero. If $c=\infty$, the previous paragraph gives convergence to zero. Both cases contradict the lower bound $\delta$.
\end{proof}

\subsection{Proof of the conditional miss rate of the Skellam rule}
\begin{proof}
Write $G_n$ for the CDF of the null distribution at candidate $n$ and $Y$ for the observation. The rule rejects $n$ if and only if $G_n(Y)\le u$. We use the standard argument for test inversion and the conservativeness of a discrete p-value \cite{neyman1937}. If $Y\sim G_n$, then $\Pr(G_n(Y)\le u)\le u$. If no lattice point satisfies $G_n(y)\le u$, the rejection event is empty. Otherwise, let $y^*$ be the largest lattice point satisfying $G_n(y)\le u$. The rejection event is then $\{Y\le y^*\}$, with probability $G_n(y^*)\le u$. Since the rejection region $\{Y\le y^*\}$ is a lower set, the rejection probabilities under the true distribution and $G_n$ differ by at most their Kolmogorov distance. Taking $n=d$ and applying Theorem~\ref{thm:uniform}, we obtain a rejection probability of at most $u+o(1)$ for the true candidate, uniformly over all $d\ge2$. The event $\{U<d\}$ implies rejection of the true candidate, giving \eqref{eq:coverage}. The argument applies to arbitrary acceptance sets and allows candidate rejection probabilities to vary nonmonotonically with $n$.
\end{proof}

Using $\le u$ as the rejection criterion with the left limit $G_n(Y-)=\Pr(Y^{\prime}<Y)$ (where $Y^{\prime}\sim G_n$) gives a different rule. The equally weighted $\pm2$ distribution has left limit zero at $-2$. This rejects its smallest support point, which has probability $1/2$, and the conditional miss rate tends to $1/2$ when $d\ge4$ and $d=o(\sqrt M)$. At that point, $G_n$ equals $1/2>u$, so the Skellam rule accepts the candidate.

For $d=3$, the observation may lie outside the support of $G_3$, and $\widetilde d\ge1$ always holds. If all three coefficients have the same sign, $\sum_vC_v^2\ge3k$. Otherwise, label the edge whose sign differs from the other two as edge $3$. The sum of squares is $3k+2(H_{12}-H_{13}-H_{23})$. Among the $k$ cells of edge $3$, at least $H_{13}+H_{23}-H_{12}$ are shared with edge $1$ or edge $2$. Thus $H_{13}+H_{23}-H_{12}\le k$, and the sum of squares is at least $k$. Hence $G_3(\widetilde d)\ge G_3(1)=1/2>u$, and the true candidate is always accepted.

\subsection{Combining the two rules using the first-round observation}\label{app:switch}
The combined rule uses the Gaussian rule's upper bound when $\widetilde d$ is at or above the switch point and the Skellam rule's upper bound otherwise. Assume that the ratio of the switch point to $\sqrt{M}$ tends to infinity as $M\to\infty$. We show that the combined rule satisfies \eqref{eq:coverage}.

\begin{proof}
Suppose that there exist $\delta>0$ and a sequence $(M,d)$ along which the combined rule satisfies $\Pr(U<d\mid F)\ge u+\delta$. Take a subsequence with $d/\sqrt M\to c\in[0,\infty]$.

For $c<\infty$, Theorem~\ref{thm:uniform} shows that the distance between the post-failure distribution of $\widetilde d-d$ and the distribution of $J_{\lambda_d}$ tends to zero. Since $\lambda_d$ is bounded, $\widetilde d=d+O_p(1)=O_p(\sqrt M)$. The probability of exceeding the switch point therefore tends to zero. The combined rule differs from the Skellam rule only on this event, so Theorem~\ref{thm:coverage} bounds the miss probability by $u+o(1)$.

For $c=\infty$, a miss implies that the rule selected from the observation rejects the true candidate $d$. The Skellam rule rejects $d$ if and only if $(\widetilde d-d)/s(d)$ is below the lower $u$-quantile of the standardized $G_d$. By the proof of Theorem~\ref{thm:uniform}, $J_{\lambda_d}/s(d)\Rightarrow\mathcal N(0,1)$, so this quantile tends to $\Phi^{-1}(u)$. The Gaussian rule rejects $d$ if and only if $(\widehat d-d)/s(d)<(t(d)-d)/s(d)$. The right-hand side also tends to $\Phi^{-1}(u)$, and $(\widehat d-\widetilde d)/s(d)\to0$ (Appendix~\ref{app:ams}). By Theorem~\ref{thm:clt}, the conditional distribution of the standardized statistic converges to the continuous distribution $\mathcal N(0,1)$. The two events for rejecting $d$ therefore differ only in a boundary region whose probability tends to zero. The combined rule's rejection event is contained in their union. Its probability is thus at most the Skellam rule's rejection probability plus $o(1)$, namely $u+o(1)$.

Both cases contradict $\delta>0$.
\end{proof}

\subsection{Extension to multiple sign counts}\label{app:multi}
Suppose each cell stores $L$ copies of the count, using independent sign hashes and a shared mapping. Conditional on $\mathcal G$, the estimates $\widehat d$ from these copies are mutually independent, each with mean $d$. Averaging independent estimates is the AMS approach to reducing variance \cite{alon1999spacecomplexity}. From $\E[\widehat d^2\mid\mathcal G]\le3d^2$ and $e^{-x}\le1-x+x^2/2$, we obtain
\begin{equation*}
\E\Bigl[\exp\Bigl(-\frac{\widehat d}{5d}\Bigr)\Bigm|\mathcal G\Bigr]\le\frac{43}{50},
\end{equation*}
and the exponential Markov inequality (Chernoff bound) applied to $U=\frac2L\sum_\ell\widehat d^{(\ell)}$ gives $\Pr(U<d\mid\mathcal G)\le(\frac{43}{50}e^{1/10})^L$. This bound holds for every deterministic mapping, hence for arbitrarily rare failure. For $L=272$, the right-hand side is below $10^{-6}$. When only an extra scalar AMS accumulator is stored, the $L$ count sums are independent given $d$, and their joint distribution is a product of binomial distributions. This permits an exact discrete candidate test with the finite-sample guarantee $\Pr(U<d\mid F)\le u$.

Repeating counts reduces absolute error and preserves the ratio of the conditional scale to the unconditional scale. Fix $d\ge4$ and $L$, and let $M\to\infty$. If we apply Gaussian inversion to the average of $L$ copies using the unconditional standard deviation $s(n)/\sqrt L$, the conditional miss rate tends to $2^{-L}\sum_{j<L/2}\binom Lj$. This equals $1/2,1/4,5/16,93/256$ for $L=1,2,4,8$, respectively. Calibration with more copies still requires the correct post-failure distribution, motivating our use of a single count and the discrete limit distribution in the main text.
